\chapter{Benchmark Methodology}

\section{Benchmark Goals}

heapx benchmarks are designed to answer two different questions.
The heap-only benchmark asks how the heap API and backend algorithms behave in isolation.
The Dijkstra benchmark asks how the backends behave inside a graph algorithm that naturally uses decrease-key.

Both benchmarks also act as consistency checks.
They compute checksums or verify monotonic extraction so performance runs can detect obvious behavioral regressions.

\section{Heap-Only Benchmark}

\texttt{tests/heap\_benchmark.c} runs all three backends over the same set of scenarios.
Items contain a key, an id, a live flag, and an optional handle.
The comparator orders by key and then id, making benchmark ordering deterministic even when keys collide.

\subsection{Scenarios}

The benchmark scenarios are:

\begin{itemize}
    \item \texttt{insert\_only}: measures insertion and minimum maintenance without draining the heap;
    \item \texttt{first\_extract}: builds a heap and extracts only the first minimum;
    \item \texttt{insert\_extract}: inserts all items and drains the heap, verifying monotonic order;
    \item \texttt{decrease\_only}: inserts handled items, decreases all priorities, and checks the new minimum;
    \item \texttt{decrease\_extract}: inserts handled items, decreases priorities, and drains the heap;
    \item \texttt{mixed\_handles}: combines handled insertion, remove, decrease-key, and extraction.
\end{itemize}

This scenario split is useful because different heaps should excel on different workloads.
For example, decrease-heavy workloads are where Fibonacci-style heaps are theoretically attractive, while insert/extract workloads often favor simpler memory layouts.

\subsection{Determinism and Output}

The benchmark uses a deterministic 64-bit linear congruential generator.
The seed can be passed from the command line or through Makefile variables.
Output can be human-readable or TSV.

\begin{lstlisting}[style=shellstyle, caption={Heap benchmark commands.}]
make benchmark-heap
make benchmark-heap N=500000 SEED=123
make benchmark-heap FORMAT=tsv N=500000 SEED=123
\end{lstlisting}

The TSV output records implementation, scenario, item count, seed, elapsed seconds, operation count, operations per second, checksum, and compiler version.

\section{TSV Comparison Workflow}

The Makefile can save benchmark TSV files and compare two runs:

\begin{lstlisting}[style=shellstyle, caption={Saving and comparing benchmark runs.}]
make benchmark-heap-tsv N=500000 SEED=123 \
  BENCHMARK_HEAP_TSV=benchmarks/before.tsv

make benchmark-heap-tsv N=500000 SEED=123 \
  BENCHMARK_HEAP_TSV=benchmarks/after.tsv

make benchmark-heap-compare \
  BASE=benchmarks/before.tsv \
  HEAD=benchmarks/after.tsv
\end{lstlisting}

Generated TSV files are ignored by git.
This keeps the repository clean while still making local performance comparisons repeatable.

\section{Dijkstra Benchmark}

\texttt{tests/dijkstra\_benchmark.c} loads a DIMACS shortest-path graph, converts it into compressed sparse row form, and runs Dijkstra's algorithm with each heap backend.

Each graph node stores:

\begin{itemize}
    \item node id;
    \item current distance;
    \item finalized flag;
    \item queued flag;
    \item heap handle.
\end{itemize}

When relaxation improves the distance of a node already in the heap, the benchmark calls \texttt{heapx\_decrease\_key()}.
When the node is not queued, it calls \texttt{heapx\_insert\_handle()}.

\section{DIMACS Parser}

The DIMACS parser is deliberately stricter than a quick benchmark parser.
It checks:

\begin{itemize}
    \item line length;
    \item problem-line uniqueness;
    \item problem type;
    \item numeric limits;
    \item arc count consistency;
    \item tail and head range;
    \item unsupported line tags;
    \item allocation overflow in graph arrays.
\end{itemize}

This matters because benchmark inputs can be large.
A permissive parser could hide errors in datasets or cause undefined behavior on oversized inputs.

\section{Cross-Backend Consistency}

After each Dijkstra run, the benchmark computes a checksum over final distances and counts reachable nodes.
All backends must agree.
This makes the benchmark more than a timing tool: it is also an end-to-end consistency check for targeted operations under a realistic workload.

\section{Interpreting Results}

Benchmark results should be interpreted carefully.
They depend on:

\begin{itemize}
    \item compiler and optimization flags;
    \item machine and cache behavior;
    \item allocator behavior;
    \item item count;
    \item graph structure;
    \item ratio of insert, decrease-key, remove, and extract-min operations.
\end{itemize}

The project therefore records enough metadata to compare runs, but avoids claiming universal performance conclusions from smoke tests.
